% =====================================================================
%  ME 306 -- sCO2 Brayton cycle project
%  Final report
%  Build:  python reports/build.py
% =====================================================================
\documentclass[11pt,a4paper]{article}

\usepackage[margin=2.5cm]{geometry}
\usepackage{amsmath,amssymb}
\usepackage{siunitx}
\usepackage{booktabs}
\usepackage{lmodern}
\usepackage[T1]{fontenc}
\usepackage{microtype}
\usepackage{graphicx}
\usepackage{tikz}
\usepackage{pgfplots}
\usepackage{listings}
\usepackage[most]{tcolorbox}
\usepackage[colorlinks=true,linkcolor=black,citecolor=black,urlcolor=blue]{hyperref}

\pgfplotsset{compat=1.17}
\usetikzlibrary{arrows.meta,positioning,calc,fit,backgrounds}
\sisetup{per-mode=symbol, detect-weight=true}

% ---------------------------------------------------------------- styles
\definecolor{cold}{RGB}{214,230,246}
\definecolor{hot}{RGB}{250,228,215}
\definecolor{neut}{RGB}{238,238,234}
\definecolor{accent}{RGB}{83,74,183}
\definecolor{warn}{RGB}{153,60,29}
\definecolor{chalk}{RGB}{47,102,78}
\definecolor{chalkbg}{RGB}{240,247,243}

\tikzset{
  blk/.style   ={draw=black!55, rounded corners=2pt, align=center,
                 font=\small, minimum height=8mm, minimum width=22mm,
                 inner sep=3pt, line width=0.4pt},
  coldblk/.style={blk, fill=cold},
  hotblk/.style ={blk, fill=hot},
  neutblk/.style={blk, fill=neut},
  ar/.style    ={-{Stealth[length=2mm,width=1.6mm]}, line width=0.5pt},
  lbl/.style   ={font=\scriptsize, inner sep=1.5pt},
}

% caveat box: every one states the choice and what it costs
\newtcolorbox{caveat}[1][]{
  enhanced, breakable, colback=warn!4, colframe=warn!60,
  boxrule=0.4pt, left=3mm, right=3mm, top=2mm, bottom=2mm,
  fonttitle=\bfseries\small, coltitle=warn!80!black,
  title={Caveat},#1}

\newtcolorbox{choicebox}[1][]{
  enhanced, breakable, colback=accent!4, colframe=accent!50,
  boxrule=0.4pt, left=3mm, right=3mm, top=2mm, bottom=2mm,
  fonttitle=\bfseries\small, coltitle=accent!80!black,
  title={Implementation},#1}

\newcommand{\code}[1]{\texttt{\small #1}}

% ------- code excerpts, pulled from source by make_code_snippets.py -------
\lstdefinestyle{pyst}{
  language=Python,
  basicstyle=\ttfamily\scriptsize,
  keywordstyle=\color{chalk!95!black}\bfseries,
  commentstyle=\color{chalk!50!black}\itshape,
  stringstyle=\color{warn!75!black},
  showstringspaces=false,
  breaklines=true, breakatwhitespace=true,
  columns=fullflexible, keepspaces=true,
  tabsize=4, aboveskip=0pt, belowskip=0pt,
}

\newtcbinputlisting{\codefile}[2][]{
  listing only, enhanced, breakable,
  colback=chalkbg, colframe=chalk!55,
  boxrule=0.4pt, left=3mm, right=3mm, top=2mm, bottom=1.5mm,
  fonttitle=\bfseries\scriptsize\ttfamily, coltitle=chalk!85!black,
  attach boxed title to top left={xshift=4mm, yshift=-2.4mm},
  boxed title style={colback=chalkbg, colframe=chalk!55, boxrule=0.4pt,
                     sharp corners, size=small},
  listing file={code/#2}, listing style=pyst, #1}

% ---------------------------------------------------------------- title
\title{\vspace{-1.5cm}\Large\bfseries
  Supercritical \ce{CO2} Brayton Cycle for a Compact SMR\\[2mm]
  \large Theory and Implementation}
\author{\normalsize ME 306 -- Advanced Thermodynamics}
\date{\normalsize Living document. Last rebuilt \today.}

\usepackage[version=4]{mhchem}



\title{\vspace{-1.5cm}\textbf{Supercritical \ce{CO2} Brayton Cycle for a
100\,MW\textsubscript{th} Compact SMR}\\[2mm]
\large Layout and recuperator trade-off under size and mass constraints}
\author{ME 306 --- Advanced Thermodynamics}
\date{}

\begin{document}
\maketitle
\vspace{-1cm}

\begin{abstract}
\noindent
A supercritical \ce{CO2} closed Brayton cycle was modelled for a
\SI{100}{\mega\watt} thermal compact SMR and four configurations were
compared at equal thermodynamic quality of heat exchange: simple
recuperated, recompression, partial cooling and intercooled. The
decision variables named in the brief were swept, candidates were
filtered against its constraints, and a Pareto front was constructed
over efficiency, mass, cost proxy and operating margin.

\textbf{The recommendation is a recompression cycle at
\SI{20}{\mega\pascal} with a split fraction of \num{0.25} and a
\SI{35}{\celsius} compressor inlet}, giving \SI{41.41}{\percent}
thermal efficiency and a second-law efficiency of \num{0.684} for
\SI{87.0}{\tonne} of power-conversion hardware ---
\SI{0.476}{\mega\watt_e\per\tonne}, with \SI{4.02}{\kelvin} of margin
above the critical temperature. Against the best feasible recuperated
design it buys \num{4.06} efficiency points for \SI{18.5}{\tonne},
paying back in \num{2.3} years on the brief's cost proxy.

Two results are worth more than the selection itself. The optimum
compressor inlet is \emph{not} at the critical point: moving from
\num{32} to \SI{35}{\celsius} makes compression about
\SI{60}{\percent} more expensive per kilogram and efficiency
nonetheless \emph{rises}, because the warmer inlet eases the
capacity-rate mismatch and the reactor supplies \SI{15}{\percent} less
heat. And intercooling loses at every intermediate pressure swept: it
does reduce compressor work, but the heat it dumps is heat the
recuperator would have recovered.

Every number in this report was produced by a script in the
repository and read from a saved result file; none was transcribed by
hand. The mass, cost and optimisation conclusions rest on the brief's
illustrative coefficients and are qualified accordingly throughout.
\end{abstract}

\tableofcontents
\newpage

% Where this document keeps the generated tables, plot data and
% schematics. The body is shared with the master document, which sits in
% a different directory and defines these differently -- so the body
% itself contains no relative paths.
\newcommand{\tbl}[1]{../documentation/latex/tables/#1}
\newcommand{\dat}[1]{../documentation/latex/data/#1}
\newcommand{\dgm}[1]{../documentation/latex/diagrams/#1}

% =====================================================================
\section{Introduction and objectives}

A small modular reactor is factory-built and shipped, so its power
conversion system has to be compact. That is the entire reason for
considering supercritical \ce{CO2} rather than steam: near its critical
point (\SI{7.38}{\mega\pascal}, \SI{30.98}{\celsius}) \ce{CO2} is dense
and compresses with very little work, which shrinks both the compressor
duty and the machinery.

The difficulty is that the benefit depends on the layout, and the
layouts that raise efficiency add hardware --- which attacks the
compactness that motivated the choice in the first place. The task is
therefore not to maximise efficiency. It is to find where the extra
hardware stops paying for itself, and to defend one configuration.

\subsection*{What was required}

\begin{center}
\small
\begin{tabular}{@{}p{5.6cm}p{8.4cm}@{}}
\toprule
Objective 1 & Maximise net thermal efficiency \\
Objective 2 & Minimise PCS mass and volume \\
Objective 3 & Minimise the cost proxy \\
Objective 4 & Maximise compressor-inlet margin above the critical
point \\
Reported, not optimised & Overall second-law efficiency \\
\midrule
\multicolumn{2}{@{}l}{\textbf{Constraints}} \\
\multicolumn{2}{@{}p{14cm}}{Compressor inlet $\ge$
\SI{32}{\celsius} and $\ge$ \SI{7.7}{\mega\pascal}; recuperator pinch
$\ge$ \SI{5}{\kelvin}; cycle pressure $\le$ \SI{20}{\mega\pascal};
turbine inlet $\le$ \SI{500}{\celsius}; split fraction $\le 0.5$.} \\
\bottomrule
\end{tabular}
\end{center}

\subsection*{Design basis}

\SI{100}{\mega\watt} thermal, indirect heater. Turbine inlet
\SI{500}{\celsius} at \SI{20}{\mega\pascal}; compressor inlet
\SI{32}{\celsius} at \SI{7.7}{\mega\pascal}; compressor isentropic
efficiency \num{0.89}, turbine \num{0.90}; recuperator effectiveness
\num{0.90} on the first pass, superseded by a \SI{5}{\kelvin} pinch
specification for all comparisons; cooling water at
\SI{25}{\celsius}; \SI{7000}{\hour} of operation a year.

\begin{caveat}[title={Scope of this report}]
This report presents the engineering investigation and its results. The
solver, its governing equations, numerical methods, software
architecture and study workflows are documented separately in
\emph{Theory and Implementation}, and are referred to here only where a
result cannot be defended without them.
\end{caveat}

% =====================================================================
\section{Methodology}

\subsection{Modelling approach}

Each cycle is a flowsheet of generic thermodynamic blocks --- turbine,
compressors, heat exchangers, coolers, splitter, mixer --- connected by
streams, and solved for the state at every point. Properties come from
CoolProp's Helmholtz-energy equation of state for \ce{CO2}. Ideal-gas
relations are not used anywhere: near the critical point $c_p$ varies by
an order of magnitude over a few kelvin, and constant-$c_p$ relations
are not approximately wrong there but qualitatively wrong.

States are carried as $(P,h)$ rather than $(P,T)$. Every block's
governing equation is an energy balance, which is an equation in
enthalpy; and in the critical region $T$ and $P$ do not uniquely fix a
state in a numerically comfortable way, while $P$ and $h$ do.

\subsection{How a number gets into this report}

No figure here was typed from a screen. Each one travels the same path:
a case file fixes the inputs, the solver converges the flowsheet, the
assessment layer turns the converged states into efficiency, mass and
exergy, and a study script writes the result to \code{results/} together
with a manifest recording the fingerprint of the source that produced
it. A generator then turns those files into the \LaTeX{} tables this
document includes. The chain is one-directional and every link is a
file on disk, so any number can be traced back to the commit-state of
the code that made it.

\begin{figure}[htbp]
\centering
\input{\dgm{pipeline}}
\caption{From a case to a printed number. The only hand-written step is
the first box.}
\label{r:fig:pipeline}
\end{figure}

\subsection{Specifying heat exchangers}

An exchanger can be specified three ways --- effectiveness, minimum
pinch, or conductance $UA$ --- and the three are not interchangeable as
a basis for comparison.

\textbf{All comparisons in this report use a common
\SI{5}{\kelvin} minimum pinch}, not a common effectiveness. Specifying
effectiveness asks each layout for the same \emph{fraction} of the
thermodynamically available heat, which is a different physical demand
in each; specifying the pinch asks each for the same \emph{quality} of
heat exchange, and lets the required size differ. Since size is one of
the four objectives, that is the only comparison that means anything.

\subsection{Sizing and cost}

The brief supplies one conversion from thermodynamics to hardware: a
printed-circuit exchanger weighs about
\SI{2.5}{\kilo\gram} per \si{\kilo\watt\per\kelvin} of $UA$. Everything
in the mass column follows from it, which is why the conductance of
every exchanger --- recuperators \emph{and} coolers --- is computed
rather than assumed.

\begin{caveat}[title={Two readings the brief does not settle}]
\textbf{The precooler.} The \SI{2.5}{\kilo\gram\per\kilo\watt\per\kelvin}
figure is given for the recuperator; the precooler is said only to
``scale with heat rejected''. Both are printed-circuit exchangers in the
same technology, so the same conversion is applied to both. This is
consequential: the precooler is the \emph{larger} exchanger in every
layout, so it carries more mass than the recuperator does.

\textbf{Volume.} A baseline volume (\SI{120}{\cubic\metre}) and mass
(\SI{60}{\tonne}) are given, but no volume-per-$UA$. Their ratio,
\SI{0.5}{\tonne\per\cubic\metre}, is used as a density. Volume therefore
carries no information that mass does not, and
\si{\mega\watt_e\per\cubic\metre} is the \si{\mega\watt_e\per\tonne}
figure rescaled. It is reported because the brief asks for it.
\end{caveat}

% =====================================================================
\section{The configurations}
\label{r:sec:configs}

All five are built from the same blocks --- compressors, a turbine, heat
exchangers, coolers, a splitter and a mixer. What distinguishes them is
only how those blocks are wired, and each successive layout is an answer
to a specific defect in the one before it.

\begin{figure}[htbp]
\centering
\input{\dgm{simple}}
\caption{\textbf{Simple Brayton.} Compress, heat, expand, reject. The
turbine exhaust still leaves at roughly \SI{390}{\celsius} and all of
that heat goes to the water, so the reactor must supply it again from
cold. Included as a reference point, not a candidate.}
\label{r:fig:simple}
\end{figure}

\begin{figure}[htbp]
\centering
\input{\dgm{recuperated}}
\caption{\textbf{Recuperated.} One exchanger moves heat from the turbine
exhaust to the compressor discharge, so the reactor starts from a
warmer state. This is the baseline every other layout is measured
against, and it recovers most of what the simple cycle wasted.}
\label{r:fig:recuperated}
\end{figure}

\begin{figure}[htbp]
\centering
\input{\dgm{recompression}}
\caption{\textbf{Recompression.} A fraction $x$ of the flow bypasses the
precooler and is compressed warm, rejoining before the
high-temperature recuperator. The low-temperature recuperator therefore
carries full flow on its hot side and only $1-x$ on its cold side ---
which is the whole point, and is explained in \S\ref{r:sec:why}.}
\label{r:fig:recompression}
\end{figure}

\begin{figure}[htbp]
\centering
\input{\dgm{partialcooling}}
\caption{\textbf{Partial cooling.} A third pressure level. The turbine
exhausts below the main compressor inlet pressure and a pre-compressor
brings it back up, so the expansion can be widened without giving up
the dense main-compressor inlet. It is the pre-compressor inlet, at
roughly \SI{6}{\mega\pascal}, that the brief's pressure constraint
turns on.}
\label{r:fig:partialcooling}
\end{figure}

\begin{figure}[htbp]
\centering
\input{\dgm{intercooled}}
\caption{\textbf{Intercooled.} Compression split in two with cooling
between the stages, so the second stage works on denser fluid. The
intercooler is the only block here that rejects heat the recuperator
could otherwise have used --- which, as \S\ref{r:sec:ic} shows, is what
decides its fate.}
\label{r:fig:intercooled}
\end{figure}

\section{Cycle configuration comparison}
\label{r:sec:layouts}

Four of the five were carried through the full comparison; the simple
cycle is omitted because without recuperation it reaches only about
\SI{23}{\percent} and is not a candidate for this duty.

All four are solved at \SI{100}{\mega\watt} thermal, \SI{500}{\celsius}
turbine inlet, water cooling, and --- importantly --- a common
\SI{5}{\kelvin} recuperator pinch rather than a common effectiveness.
That choice is what makes the comparison meaningful: it asks each layout
for the same \emph{quality} of heat exchange and lets the required size
differ, so the hardware each one needs is an output of the comparison
rather than an assumption built into it.

\begin{table}[htbp]
\centering
\input{\tbl{layout_comparison}}
\caption{The four configurations at equal pinch. Conductances are in
\si{\kilo\watt\per\kelvin}.}
\label{r:tab:layouts}
\end{table}

\begin{table}[htbp]
\centering
\input{\tbl{total_conductance}}
\caption{Total exchanger conductance --- recuperators \emph{and}
coolers. The last column is the ratio to the recuperated baseline.}
\label{r:tab:conductance}
\end{table}

\subsection{Recompression wins, and counting coolers changes by how much}

Recompression gains \num{5.71} efficiency points over the recuperated
baseline. The hardware cost of those points depends entirely on what is
counted.

Counting \emph{recuperators only}, recompression needs
\SI{9936}{\kilo\watt\per\kelvin} against \SI{3967}{} --- a factor of
\num{2.51}. Counting every exchanger, it needs
\SI{17598}{\kilo\watt\per\kelvin} against \SI{12386}{} --- a factor of
\num{1.42}.

The difference is not an accounting convenience. Recompression makes
the precooler \emph{smaller}, for two reasons: it rejects less heat per
megawatt, because the recuperators recover more of it; and the
precooler carries only the unsplit fraction of the flow, because the
recompressed branch bypasses it entirely. Charging recompression for
the recuperator it added while ignoring the cooler it saved overstates
its cost by three-quarters. Every comparison in this report is
therefore on total conductance.

The same correction moves partial cooling from \num{1.35} to
\num{1.07}.

\subsection{Why recompression works}
\label{r:sec:why}

A recuperator fails when the two streams have different heat capacity
rates $\dot m c_p$. The cold stream, at \SI{20}{\mega\pascal} and
leaving the compressor near \SI{60}{\celsius}, sits close to the
pseudo-critical line where $c_p$ is large; the hot stream at
\SI{7.7}{\mega\pascal} does not. The cold side therefore needs far more
heat per kelvin than the hot side can give up, the profiles converge in
the middle rather than at an end, and the exchanger pinches internally
where no amount of area helps.

Splitting the flow fixes it directly. Sending a fraction $x$ around the
precooler and compressing it separately reduces the mass flow on the
cold side of the low-temperature recuperator, which reduces its
capacity rate, which lets the profiles stay apart. The split is not an
efficiency trick; it is a capacity-rate correction.

\subsection{Intercooling loses, and the reason is not the pressure}
\label{r:sec:ic}

Intercooling has a free variable --- the intermediate pressure --- and
it was swept rather than assumed.

\begin{table}[htbp]
\centering
\input{\tbl{intercooling}}
\caption{Intercooling against the plain recuperated cycle
(\SI{36.541}{\percent}) at every intermediate pressure swept. The lower
end approaches the compressor inlet, where the intercooler effectively
switches off.}
\label{r:tab:intercooling}
\end{table}

Three findings, in increasing order of interest.

\textbf{The intercooler does what it is supposed to do.} Back-work ratio
falls from \num{0.176} without it to \num{0.160} at its best point: the
compressors really do take back less of the turbine's work.

\textbf{The pressure that minimises compressor work is not the pressure
that maximises efficiency.} Back-work ratio bottoms out at
\SI{10}{\mega\pascal}, while efficiency is still improving at
\SI{7.9}{}. Optimising the thing intercooling is \emph{for} would select
the wrong pressure.

\textbf{There is no interior optimum at all.} Efficiency rises
monotonically as the intermediate pressure falls towards the compressor
inlet --- that is, as the intercooler is progressively switched off. The
best intercooler for this cycle is no intercooler.

The reason is the second, unintended effect. Heat dumped in the
intercooler leaves at a temperature the recuperator could have used, so
the reactor must supply it again. In a gas turbine, with a back-work
ratio near \num{0.5}, the compression saving dominates and intercooling
is standard practice. In \ce{CO2} compression is already cheap, so there
is little to save and the recuperator penalty wins.

The classical geometric-mean rule
($\sqrt{7.7 \times 20} = \SI{12.41}{\mega\pascal}$) is an ideal-gas
result and sits \num{0.23} points below the best point swept --- a
reminder that near the critical point the textbook rules are starting
points, not answers.

\subsection{Partial cooling: performance and feasibility are separate}

Partial cooling adds a third pressure level, the turbine exhaust, which
the brief does not constrain. It too was swept.

\begin{table}[htbp]
\centering
\input{\tbl{partial_cooling}}
\caption{Partial cooling's turbine exhaust pressure, split held at
\num{0.30}. The last column marks whether the pre-compressor inlet
satisfies the brief's \SI{7.7}{\mega\pascal} limit under its literal
reading.}
\label{r:tab:partialcooling}
\end{table}

\textbf{On performance:} \SI{6.0}{\mega\pascal} is an interior optimum
within the range swept. Re-sweeping the split at that pressure gives
\SI{41.9}{\percent} at $x = 0.40$, against recompression's
\SI{44.9}{\percent} at the same split. Partial cooling does not
overtake recompression at any point investigated.

Two pressures do not solve, for a physical rather than a numerical
reason: as the exhaust approaches the main compressor inlet, the
pre-compressor has almost no work left to do, the second precooler's
duty collapses, and its coolant temperature profile crosses.

\textbf{On feasibility, which is a separate question:} none of the five
solvable exhaust pressures places the pre-compressor inlet at or above
\SI{7.7}{\mega\pascal}. Under the literal reading of the brief's
constraint, the entire partial-cooling family is inadmissible whatever
its efficiency, and the \SI{41.9}{\percent} figure above is therefore
\emph{not} a feasible design under the stated requirement. This is
discussed in \S\ref{r:sec:ambiguity}.

% =====================================================================
\section{Component assessment}

\subsection{Turbine and compressors}

At \SI{500}{\celsius} and \SI{20}{\mega\pascal} expanding to
\SI{7.7}{\mega\pascal} at \num{0.90} isentropic efficiency, the turbine
exit lands near \SI{391}{\celsius} --- which is what sets the
hot-side recuperator inlet and therefore how much heat can be recovered
at all.

The compressors are where sCO\textsubscript{2} earns its reputation. The
back-work ratio across the four layouts runs from \num{0.16} to
\num{0.28}, against roughly \num{0.5} for a gas turbine and \num{0.01}
for a steam plant. Compression near the critical point is cheap because
the fluid is dense and $w = \int v\,\mathrm{d}P$ is small --- the
compressibility factor at the inlet is about \num{0.22}.

That cheapness is also why intercooling fails here: there is simply not
much compression cost available to save.

\subsection{Recuperators}

The recuperator is the component the whole layout question is about. At
a \SI{5}{\kelvin} pinch the recompression layout needs
\SI{9936}{\kilo\watt\per\kelvin} split across two units, against
\SI{3967}{} for the single recuperated unit --- and that extra
conductance is the mass the efficiency is bought with.

Its minimum temperature difference is not at either end. Because of the
capacity-rate mismatch the profiles converge internally, so the pinch
must be found by scanning the whole profile; checking the endpoints
reports a comfortable margin where the real one may be zero.

\subsection{Precooler and heat rejection}

\begin{table}[htbp]
\centering
\input{\tbl{coolers}}
\caption{Every cooler in every layout, sized against
\SI{25}{\celsius} water with a \SI{10}{\kelvin} rise. Conductances in
\si{\kilo\watt\per\kelvin}. A dash marks a quantity that does not exist
for that row.}
\label{r:tab:coolers}
\end{table}

The precooler is the largest exchanger in every layout, not the
recuperator. It rejects more heat, against a much smaller temperature
difference, and the combination is what drives its conductance. Any
mass estimate that counts recuperators and omits coolers is wrong by
more than a factor of two.

\begin{table}[htbp]
\centering
\input{\tbl{water_sweep}}
\caption{The baseline precooler as cooling-water flow is increased.}
\label{r:tab:water}
\end{table}

Adding water shrinks the exchanger --- \SI{12450}{} to
\SI{5933}{\kilo\watt\per\kelvin} --- until the pinch reaches
\SI{7.00}{\kelvin} and stops. That floor is the \SI{32}{\celsius}
\ce{CO2} outlet minus the \SI{25}{\celsius} water inlet, and no flow
opens it further. Meanwhile $\eta_{th}$ does not move at all: only
$\eta_{net}$ changes, and it \emph{falls}, because pump work is a cost
with no efficiency benefit to trade against. Cooling water is a siting
constraint and a hardware trade, not an efficiency lever.

% =====================================================================
\section{Second-law analysis}

\begin{table}[htbp]
\centering
\input{\tbl{exergy_balance}}
\caption{Exergy balance by layout, \si{\mega\watt}, against a
\SI{25}{\celsius} dead state.}
\label{r:tab:exergy}
\end{table}

For the recommended recompression design, the four largest
irreversibilities are the high-temperature recuperator
(\SI{6.09}{\mega\watt}, \SI{45.3}{\percent} of the accounted total), the
turbine (\SI{2.90}{}, \SI{21.6}{\percent}), the precooler
(\SI{1.82}{}, \SI{13.5}{\percent}) and the low-temperature recuperator
(\SI{0.90}{}, \SI{6.7}{\percent}).

\textbf{The recuperator dominates in every layout.} That is the
capacity-rate mismatch of \S\ref{r:sec:layouts} seen a second way: heat
crossing a finite temperature difference destroys availability, and the
mismatch is precisely what forces that difference to stay open along
most of the exchanger. The first-law and second-law analyses are not
two findings; they are one finding in two currencies.

\begin{caveat}[title={The totals are incomplete, and say so}]
The heater's destruction cannot be computed: the model takes the
reactor as given and has no heat-source boundary temperature, so the
entropy generated transferring heat into the cycle is unknown. Every
total above excludes it and is labelled incomplete rather than
presented as a total.

This is not a small omission. With a plausible \SI{650}{\celsius}
source the heater would destroy about \SI{13.8}{\mega\watt} --- more
than the recuperator --- and would be the largest irreversibility in the
cycle. The layouts remain comparable because the same term is missing
from each, but no absolute figure here is a complete accounting.
\end{caveat}

The overall second-law efficiency of the recommended design is
\textbf{\num{0.684}}, against a Carnot efficiency between the assumed
source and sink temperatures.

% =====================================================================
\section{Parametric studies and optimisation}

\subsection{Pressure}

\begin{figure}[htbp]
\centering
\begin{tikzpicture}
\begin{axis}[width=11.5cm, height=5.4cm,
  xlabel={Maximum cycle pressure [\si{\mega\pascal}]},
  ylabel={$\eta_{th}$ [\si{\percent}]},
  xmin=13.6, xmax=20.4, ymin=29, ymax=44,
  legend pos=south east, legend style={font=\small,draw=black!30},
  grid=both, grid style={black!8},
  tick label style={font=\small}, label style={font=\small}]
\addplot[accent, thick, mark=*, mark size=1.4pt]
  table[x=P_high,y=eta_recup] {\dat{eta_vs_rp.dat}};
\addlegendentry{recuperated}
\addplot[warn, thick, mark=square*, mark size=1.4pt]
  table[x=P_high,y=eta_recomp] {\dat{eta_vs_rp.dat}};
\addlegendentry{recompression}
\end{axis}
\end{tikzpicture}
\caption{Efficiency against maximum pressure at a fixed
\SI{7.7}{\mega\pascal} inlet, so the axis is also pressure ratio
(\numrange{1.82}{2.60}).}
\label{r:fig:rp}
\end{figure}

Both curves rise all the way to the \SI{20}{\mega\pascal} limit and are
still rising there. This is the opposite of the air-standard case,
where a real cycle at fixed turbine inlet temperature has an interior
optimum because past some ratio the compressor eats work faster than
the turbine gains it. Here the back-work ratio is low enough that the
crossover has not happened by \SI{20}{\mega\pascal}: \textbf{the
optimum is the constraint, not a property of the cycle.} Raise the
pressure limit and the answer moves.

\subsection{Split fraction}

\begin{figure}[htbp]
\centering
\begin{tikzpicture}
\begin{axis}[width=11.5cm, height=5.4cm,
  xlabel={Split fraction $x$ (fraction bypassing the precooler)},
  ylabel={$\eta_{th}$ [\si{\percent}]},
  xmin=0.07, xmax=0.48, ymin=36.5, ymax=45.5,
  legend pos=north west, legend style={font=\small,draw=black!30},
  grid=both, grid style={black!8},
  tick label style={font=\small}, label style={font=\small}]
\addplot[warn, thick, mark=square*, mark size=1.4pt]
  table[x=x,y=eta_recomp] {\dat{eta_vs_x.dat}};
\addlegendentry{recompression}
\addplot[accent, thick, mark=*, mark size=1.4pt]
  table[x=x,y=eta_partial] {\dat{eta_vs_x.dat}};
\addlegendentry{partial cooling}
\end{axis}
\end{tikzpicture}
\caption{Efficiency against split fraction at \SI{20}{\mega\pascal}.}
\label{r:fig:x}
\end{figure}

Recompression has a genuine interior optimum near $x = 0.40$ and falls
away beyond it. An interior optimum is the signature of two effects of
opposite sign, and both can be named: sending more flow around the
precooler relieves the capacity-rate mismatch, which is pure gain, until
too little flow is left passing through the precooler to absorb the heat
that must leave. Then the recuperator starts starving and the curve
turns over.

\subsection{Compressor inlet state}

This is the least obvious result in the project.

The brief's floor is \SI{32}{\celsius}, and that is where the fluid is
densest at \SI{7.7}{\mega\pascal}, so compression should be cheapest
there --- and it is, by a wide margin. Moving the inlet from \num{32} to
\SI{35}{\celsius} makes compression roughly \SI{60}{\percent} more
expensive per kilogram (\SIrange{20.5}{33.4}{\kilo\joule\per\kilogram}).

And the efficiency goes \emph{up}: \SIrange{36.541}{37.355}{\percent}
for the recuperated layout.

The reason is the recuperator again. A warmer, less dense compressor
inlet produces a warmer recuperator cold inlet, which eases the
capacity-rate mismatch, which lets the recuperator recover more heat ---
so the reactor has to supply appreciably less
(\SIrange{262}{222}{\kilo\joule\per\kilogram}, about \SI{15}{\percent}).
Compression got more expensive and the heat bill fell by more.

It also buys margin: \SI{1.02}{\kelvin} above the critical temperature
at the constraint becomes \SI{4.02}{\kelvin} at \SI{35}{\celsius}. A
design better on efficiency \emph{and} better on safety margin, for
some additional exchanger, is the rare case where two objectives stop
competing --- and it would be missed entirely by anyone who assumed the
optimum sits at the critical point because that is where the density is.

\subsection{The feasible set and the Pareto front}

\num{43} candidates were generated by sweeping maximum pressure,
compressor inlet state, split fraction and layout. \num{36} are
feasible under the brief's constraints; \num{7} are rejected. Of the
feasible set, \num{5} are non-dominated.

\begin{figure}[htbp]
\centering
\begin{tikzpicture}
\begin{axis}[width=11.8cm, height=6.4cm,
  xlabel={PCS mass [\si{\tonne}]}, ylabel={$\eta_{th}$ [\si{\percent}]},
  xmin=64, xmax=108, ymin=31, ymax=44.5,
  legend pos=south east, legend style={font=\small,draw=black!30},
  grid=both, grid style={black!8},
  tick label style={font=\small}, label style={font=\small}]
\addplot[only marks, mark=x, mark size=2.6pt, black!25]
  table[x=mass,y=eta] {\dat{pareto_dominated.dat}};
\addlegendentry{dominated}
\addplot[only marks, mark=o, mark size=2.4pt, warn!70]
  table[x=mass,y=eta] {\dat{pareto_infeasible.dat}};
\addlegendentry{infeasible}
\addplot[only marks, mark=*, mark size=2.6pt, accent]
  table[x=mass,y=eta] {\dat{pareto_front.dat}};
\addlegendentry{frontier}
\addplot[only marks, mark=star, mark size=4.5pt, chalk, thick]
  table[x=mass,y=eta] {\dat{pareto_knee.dat}};
\addlegendentry{selected}
\end{axis}
\end{tikzpicture}
\caption{Efficiency against mass, all \num{43} candidates. Several
infeasible designs sit \emph{above} the front --- the constraint filter
working, not a plotting error. Some dominated crosses also sit above
and left of a front point: dominance is judged on all four objectives,
and those designs lose on cost or margin.}
\label{r:fig:paretomass}
\end{figure}

\begin{figure}[htbp]
\centering
\begin{tikzpicture}
\begin{axis}[width=11.8cm, height=5.8cm,
  xlabel={$\eta_{th}$ [\si{\percent}]},
  ylabel={Specific power [MW$_e$ per tonne]},
  xmin=31, xmax=44.5, ymin=0.37, ymax=0.57,
  legend pos=north east, legend style={font=\small,draw=black!30},
  grid=both, grid style={black!8},
  tick label style={font=\small}, label style={font=\small}]
\addplot[only marks, mark=x, mark size=2.6pt, black!25]
  table[x=eta,y=mwe_per_t] {\dat{pareto_dominated.dat}};
\addlegendentry{dominated}
\addplot[only marks, mark=o, mark size=2.4pt, warn!70]
  table[x=eta,y=mwe_per_t] {\dat{pareto_infeasible.dat}};
\addlegendentry{infeasible}
\addplot[only marks, mark=*, mark size=2.6pt, accent]
  table[x=eta,y=mwe_per_t] {\dat{pareto_front.dat}};
\addlegendentry{frontier}
\addplot[only marks, mark=star, mark size=4.5pt, chalk, thick]
  table[x=eta,y=mwe_per_t] {\dat{pareto_knee.dat}};
\addlegendentry{selected}
\end{axis}
\end{tikzpicture}
\caption{The same designs as specific power against efficiency. The
frontier slopes \emph{down}.}
\label{r:fig:paretopower}
\end{figure}

The two Pareto plots must be read together. In
Figure~\ref{r:fig:paretomass} the frontier slopes up: every efficiency
point costs tonnes. In Figure~\ref{r:fig:paretopower} the same frontier
slopes down, because specific power is the ratio of two objectives and
exchanger mass grows faster than efficiency --- so it is maximised by
the \emph{least} efficient design on the front. A reader shown only the
second would conclude the simplest layout is best. Specific power is
reported, as the brief asks, and is not used as an objective.

\begin{table}[htbp]
\centering
\input{\tbl{pareto_front}}
\caption{The non-dominated designs.}
\label{r:tab:front}
\end{table}

\textbf{Seven rejected designs.} One fails on a recuperator pinch of
\SI{4.98}{\kelvin} against the \SI{5.00}{} limit --- not a bad design,
a design whose exchanger is fractionally too small. The other six are
the partial-cooling family, rejected on inlet pressure
(\S\ref{r:sec:ambiguity}). Notably the highest-efficiency design found
anywhere in the sweep is among the rejected.

% =====================================================================
\section{Recommended design}
\label{r:sec:recommendation}

\begin{table}[htbp]
\centering
\input{\tbl{performance}}
\caption{Recompression, \SI{20}{\mega\pascal}, $x = 0.25$,
\SI{35}{\celsius} compressor inlet. Generated from
\code{results/selected\_design/}.}
\label{r:tab:selected}
\end{table}

\begin{figure}[htbp]
\centering
\begin{tikzpicture}
\begin{axis}[width=12cm, height=8cm,
  xlabel={Specific entropy [\si{\kilo\joule\per\kilogram\per\kelvin}]},
  ylabel={Temperature [\si{\celsius}]},
  xmin=1.54, xmax=2.80, ymin=0, ymax=560,
  legend pos=south east, legend style={font=\small,draw=black!30},
  grid=both, grid style={black!8},
  nodes near coords, point meta=explicit symbolic,
  every node near coord/.append style={font=\scriptsize, anchor=south east,
                                       inner sep=1.5pt, black!75},
  tick label style={font=\small}, label style={font=\small}]
\addplot[accent, thick, mark=*, mark size=1.5pt]
  table[x=s,y=T,meta=name] {\dat{ts_path.dat}};
\addlegendentry{main loop}
\addplot[warn, thick, dashed, mark=square*, mark size=1.5pt]
  table[x=s,y=T,meta=name] {\dat{ts_branch.dat}};
\addlegendentry{recompression branch}
\end{axis}
\end{tikzpicture}
\caption{$T$--$s$ diagram of the recommended cycle. $1\to2$ main
compressor; $2\to3\to4\to5$ the two recuperators and the mixing point;
$5\to6$ reactor; $6\to7$ turbine; $7\to8\to9$ the recuperators giving
the heat back; $9\to1$ precooler. The dashed branch is the quarter of
the flow that skips the precooler, compressed from state~9 to~10 and
rejoining at~4.}
\label{r:fig:ts}
\end{figure}

Two features of the diagram carry the whole argument. \textbf{The loop
is thin}: the two recuperator legs run close together over most of
their length, and that narrow gap is the pinch, the capacity-rate
mismatch and the largest irreversibility in the cycle, all being the
same thing seen in temperature--entropy coordinates. \textbf{The
precooler leg is short}: almost all the heat leaving the hot stream is
handed to the cold stream rather than to the water, which is the
recuperator earning its tonnes.

\subsection{What was sacrificed, and to what}

The brief asks what the selection gives up relative to each role's
individually optimal design. It gives up something to all of them.

\begin{center}
\small
\begin{tabular}{@{}p{3.4cm}p{4.4cm}p{6.2cm}@{}}
\toprule
Role's optimum & Would have chosen & What the selection sacrifices \\
\midrule
R1 turbine & highest pressure ratio & nothing --- both agree on
\SI{20}{\mega\pascal}, the constraint \\[2pt]
R2 compressors & inlet at \SI{32}{\celsius}, where compression is
cheapest & \SI{60}{\percent} more compression work per
\si{\kilo\gram}, bought back on the heat bill \\[2pt]
R3 recuperators & largest affordable $UA$ & \num{1.4} efficiency
points against the $x = 0.35$ design, for \SI{17}{\tonne} \\[2pt]
R4 precooler & most cooling water & pump work and exchanger size
traded; water flow is sized to a \SI{10}{\kelvin} rise \\[2pt]
R5 compactness & the recuperated cycle, best
\si{\mega\watt_e\per\tonne} on the front &
\SI{18.5}{\tonne} and \num{0.070}
\si{\mega\watt_e\per\tonne}, for \num{4.06} efficiency points \\
\bottomrule
\end{tabular}
\end{center}

\subsection{Why this point and not another on the front}

The selection rule is a knee point: normalise each objective across the
front, construct the ideal corner where all four take their best value
at once, and take the design closest to it.

\textbf{This weights all four objectives equally, which nobody asked
for.} It is not neutral merely because no weights were typed. It is a
defensible starting point for an argument, and the argument is this:

\begin{itemize}\itemsep3pt
\item It is \emph{not} the most efficient design on the front. Two
      recompression designs beat it by up to \num{1.4} points, and both
      are heavier --- the knee gives up those points for roughly
      \SI{17}{\tonne}.
\item It is \emph{not} the most compact. The recuperated cycle at
      \SI{68.5}{\tonne} is lighter and has the best specific power on
      the front. It is also four efficiency points worse, and the
      hardware that closes that gap pays for itself in a little over
      two years against a plant meant to run for thirty.
\item \textbf{The payback argument is what decides it.} At the brief's
      coefficients, \num{4.06} extra efficiency points are worth
      \SI{1.99}{\text{US\$M}} a year against
      \SI{4.6}{\text{US\$M}} of extra hardware. That conclusion
      survives being wrong about the price of steel by a factor of two,
      which is the only reason it is worth stating.
\end{itemize}

Weight efficiency more heavily and the selection moves up the front to
the \SI{18}{\mega\pascal}, $x = 0.35$ design; weight mass more heavily
and it moves down to the recuperated cycle. \textbf{The front is the
finding. The point on it is a recommendation.}

% =====================================================================
\section{Validation, assumptions and limitations}

\subsection{What the model was checked against}

\begin{center}
\small
\begin{tabular}{@{}p{3.0cm}p{2.0cm}p{8.6cm}@{}}
\toprule
What & Status & Evidence \\
\midrule
Solver numerics & Validated & Air-standard Brayton against the closed
form, $4.6\times10^{-16}$ \\[2pt]
\ce{CO2} properties & Validated & Dostal's table to
\SI{0.002}{\percent}; three entropy disagreements traced to errors in
his table \\[2pt]
Whole cycle & Validated & Kim (2018): \SI{44.669}{\percent} against his
\SI{44.670}{}, \emph{and} his optimum found independently in two
directions \\[2pt]
Energy closure & Verified & First law closes to $2\times10^{-14}$ \\[2pt]
Exergy & Tested & 16 internal checks; no external reference available \\[2pt]
Numerical settings & Verified & Largest deviation
$9.6\times10^{-6}$ points across every tolerance swept \\[2pt]
Mass, cost, MOO & \textbf{Unvalidated} & The brief's illustrative
coefficients, applied consistently. Nothing external exists to check
them against \\
\bottomrule
\end{tabular}
\end{center}

Reproducing a published efficiency is weaker evidence than it looks ---
with enough free parameters any model can hit one number. What makes
the Kim comparison worth something is that the solver finds his
\emph{optimum}, in split fraction and pressure ratio independently,
without being told where it is. An optimum is a property of the whole
response surface, so matching it means the model's derivatives are
right and not just its value.

\textbf{The recommendation rests partly on the unvalidated row.}
Efficiency and margin are validated quantities; mass, cost, payback and
the selection itself inherit the status of the brief's coefficients.
Nothing in this report should be read as claiming otherwise.

\subsection{Numerical convergence}

Every tolerance and discretisation setting was swept across a decade or
more. The largest efficiency deviation was
$9.6\times10^{-6}$ percentage points against a $10^{-3}$ threshold for
the fourth significant figure, so quoting \SI{41.41}{\percent} is
supported.

Two findings are worth recording. A loose tear tolerance
($10^{-6}$) is rejected by the \emph{exergy} balance, residual
$5.3\times10^{-6}$, rather than by the first-law closure --- the
second-law check is the more sensitive convergence detector of the two.
And efficiency and conductance do not converge at the same rate:
exchanger resolution moves $UA$ by \SI{0.041}{\percent} while leaving
efficiency unchanged to $10^{-8}$.

\textbf{That second finding matters because mass is an objective}, so
it was propagated rather than noted. Across \numrange{25}{400} nodes
the recommended design's mass varies by \SI{0.0124}{\tonne}, against a
smallest gap between adjacent front designs of \SI{4.91}{\tonne} --- a
factor of \num{396}. The entire optimisation was then re-run at both
resolution extremes: \num{43} candidates, \num{36} feasible, and
\textbf{the same design selected at \num{25}, \num{100} and \num{400}
nodes}. Resolution does not touch the recommendation.

\subsection{The constraint that admits two readings}
\label{r:sec:ambiguity}

The brief states: \emph{compressor inlet $\ge$ \SI{32}{\celsius} and
$\ge$ \SI{7.7}{\mega\pascal} (safe distance from the two-phase
region)}. Two readings are available.

\textbf{Literal:} every compressor inlet meets both numbers.
\textbf{By intent:} every compressor inlet is a safe distance from the
two-phase region --- which, above the critical temperature, it is at any
pressure, because there is no saturation line up there to be near.

Partial cooling is exactly the case that separates them. Its
pre-compressor sits at about \SI{6}{\mega\pascal} and
\SI{32}{\celsius}: below the stated pressure, above the critical
temperature, and therefore nowhere near a phase boundary.

\textbf{This report adopts the literal reading}, because a number in a
requirement is a number, and it is the reading under which the
recommendation above holds. The consequence is explicit: the entire
partial-cooling family is excluded on a pressure limit rather than on
performance. Under the alternative reading, six further designs re-enter
the comparison and the recommendation would have to be re-examined.
This should be settled with the course staff rather than decided here.

\subsection{What is not modelled}

\textbf{Would change the ranking:} off-design and part-load behaviour
--- every efficiency here is a design-point number; turbomachinery mass
scaling with flow, held constant across layouts; the heat-source
boundary temperature, whose absence makes every exergy total partial.

\textbf{Would change the numbers, not the ranking:} pressure drops are
stated rather than derived from geometry; component efficiencies are
constant rather than functions of operating point; the cost proxy
contains no fuel, O\&M, financing or civil works.

\textbf{Open:} partial cooling was swept in one dimension at a time
rather than jointly over exhaust pressure and split fraction.

% =====================================================================
\section{Conclusions}

\begin{enumerate}\itemsep4pt
\item \textbf{Recompression at \SI{20}{\mega\pascal}, $x = 0.25$, with a
      \SI{35}{\celsius} compressor inlet} is recommended:
      \SI{41.41}{\percent} thermal efficiency, second-law efficiency
      \num{0.684}, \SI{87.0}{\tonne},
      \SI{0.476}{\mega\watt_e\per\tonne}, \SI{4.02}{\kelvin} of margin,
      \num{2.3}-year payback against the best feasible recuperated
      design.
\item \textbf{Compare layouts on total conductance, not recuperators.}
      Recompression makes the precooler smaller; counting recuperators
      alone overstates its hardware cost by three-quarters
      (\num{2.51}$\times$ against \num{1.42}$\times$).
\item \textbf{The compressor-inlet optimum is not at the critical
      point.} Warmer costs \SI{60}{\percent} more to compress and saves
      more on heat, and buys margin as well --- the one place in this
      study where two objectives stop competing.
\item \textbf{Intercooling does not pay in sCO\textsubscript{2}.} It
      reduces compressor work as intended and still loses, at every
      intermediate pressure, because the heat it dumps is heat the
      recuperator would have recovered.
\item \textbf{Partial cooling is competitive on performance and
      inadmissible as specified.} Those are separate conclusions and
      the second rests on a reading of the brief, not on physics.
\item \textbf{The pressure optimum is the constraint.} Both layouts are
      still improving at \SI{20}{\mega\pascal}; raising the limit moves
      the answer.
\item \textbf{The front is the finding, the point on it is a choice.}
      Equal weighting is a judgement, and shifting it moves the
      selection in a way anyone can inspect.
\end{enumerate}

% =====================================================================
\appendix
\section{Reference results}
\label{r:app:results}

Every table in this appendix is generated from the saved study outputs
by \code{python -m tools.make\_tables} and regenerates whenever the
report is rebuilt. None is maintained by hand.

\subsection{State points of the recommended cycle}

\begin{table}[htbp]
\centering
\input{\tbl{state_points}}
\caption{From \code{results/selected\_design/state\_points.csv}.}
\end{table}

\subsection{Numerical sensitivity}

\begin{table}[htbp]
\centering
\input{\tbl{tolerances}}
\caption{From \code{results/tolerance\_sensitivity/sweep.csv}. Rows
that did not solve are omitted; the tear tolerance of $10^{-6}$ is one
of them.}
\end{table}

\subsection{Provenance}
\label{r:sec:prov}

\begin{center}
\small
\begin{tabular}{@{}p{3.6cm}p{5.4cm}p{5.0cm}@{}}
\toprule
Result & Study & Saved to \\
\midrule
Tables \ref{r:tab:layouts}, \ref{r:tab:conductance} &
\code{compare\_layouts} & \code{results/layout\_comparison/} \\
Figures \ref{r:fig:rp}, \ref{r:fig:x} & \code{sweeps} &
\code{results/parameter\_sweeps/} \\
Table \ref{r:tab:partialcooling} & \code{partial\_cooling} &
\code{results/partial\_cooling/} \\
Table \ref{r:tab:intercooling} & \code{intercooling} &
\code{results/intercooling/} \\
Tables \ref{r:tab:coolers}, \ref{r:tab:water} & \code{cooling} &
\code{results/cooling/} \\
Table \ref{r:tab:exergy} & \code{exergy} & \code{results/exergy/} \\
Figures \ref{r:fig:paretomass}, \ref{r:fig:paretopower},
Table \ref{r:tab:front} & \code{tradeoff} &
\code{results/multi\_objective/} \\
Table \ref{r:tab:selected}, Figure \ref{r:fig:ts} &
\code{selected\_design} & \code{results/selected\_design/} \\
\bottomrule
\end{tabular}
\end{center}

Each directory carries a \code{manifest.json} recording the command,
the grids swept, the environment and a SHA-256 fingerprint of the
source that produced it.



\end{document}
